% TOP_PROBLEM: {"id": 3255, "title": "The Two Color Conjecture", "category_id": 3, "category": {"id": 3, "name": "graph_theory", "display_name": "Graph Theory", "description": "Problems involving graphs, networks, and their properties.", "slug": "graph-theory", "order_index": 3, "created_at": "2026-07-31T15:26:25.671Z"}, "status": "open", "problem_number": "OPG-169", "set_id": 12}
% FIRST_POSTED: 2026-10-01
% ENTRY_KIND: statement_only
% UnsolvedMath ID 3255; source catalogue code OPG-169
% !TeX program = lualatex
% Definition and sources only; this file is not an LLM solution attempt.
\documentclass[11pt,a4paper]{article}
\usepackage[margin=25mm]{geometry}
\usepackage{fontspec}
\setmainfont{lmroman10-regular.otf}[BoldFont=lmroman10-bold.otf,
 ItalicFont=lmroman10-italic.otf,BoldItalicFont=lmroman10-bolditalic.otf]
\usepackage{amsmath,amssymb,mathtools}
\usepackage{unicode-math}
\setmathfont{latinmodern-math.otf}
\AtBeginDocument{\let\setminus\smallsetminus}
\usepackage{xurl,hyperref}
\hypersetup{colorlinks=true,urlcolor=blue}
\setcounter{secnumdepth}{0}
\setlength{\parindent}{0pt}
\setlength{\parskip}{5pt}
\setlength{\emergencystretch}{3em}
\begin{document}
\section*{The Two Color Conjecture}\label{3255}
{\small UnsolvedMath ID 3255 · OPG-169 · Graph Theory · OpenGarden}
\subsection{Definitions and mathematical statement}
An oriented graph $D=(V,A)$ is obtained from a finite simple undirected
graph by assigning one direction to each edge. In particular it has no
loops, no repeated arcs, and no pair of opposite arcs. Assume its
underlying undirected graph is planar. A directed cycle is a cyclic
sequence of distinct vertices $v_1,\ldots,v_m$, with $m\ge3$, such that
$v_i v_{i+1}\in A$ for all indices read modulo $m$.

For $X\subseteq V$, the induced subdigraph $D[X]$ retains all arcs whose
ends lie in $X$. It is acyclic if it has no directed cycle. The
Neumann-Lara two-colour conjecture asks whether there is always a
partition
\[
                  V=X_1\sqcup X_2
                  \quad\text{with }D[X_1],D[X_2]\text{ acyclic}.
\]
Empty parts are allowed. Equivalently, there should be a map
$c:V\to\{1,2\}$ such that every directed cycle contains both colours.
The least number of acyclic induced colour classes is called the
dichromatic number, so the requested bound is at most two.

This is a directed condition. A colour class may contain an undirected
cycle if its orientation is not cyclic, and its vertices need not form
an independent set. Conversely, the prohibition of opposite arcs in an
oriented graph is material: a pair of opposite arcs would itself be a
directed cycle of length two. The assertion covers every orientation of
every finite simple planar graph, without any assumption about the
lengths of its directed cycles.
\subsection{Short English statement}
Can the vertices of every planar oriented graph be coloured with two colours so that no directed cycle is monochromatic?
\subsection{Sources}
\begin{itemize}
\item[S1] ULAM.AI and the UnsolvedMath Contributors, supplied problems.json, record 3255 (OPG-169), OpenGarden collection. Catalogue identity and source transcription; CC BY 4.0.
\url{https://huggingface.co/datasets/ulamai/UnsolvedMath}.
\item[S2] Open Problem Garden, The Two Color Conjecture. Original problem and discussion; the mathematical formulation below is expanded from the supplied source record.
\url{http://www.openproblemgarden.org/op/partitioning_planar_digraphs}.
\end{itemize}
\end{document}
